% Generated draft: requires mathematical and source review.
\documentclass[11pt]{article}
\usepackage[margin=25mm]{geometry}
\usepackage{fontspec}
\setmainfont{Latin Modern Roman}
\usepackage{amsmath,amssymb,mathtools,unicode-math}
\setmathfont{Latin Modern Math}
\usepackage{enumitem,needspace,xurl,hyperref}
\setlength{\parindent}{0pt}
\setlength{\parskip}{5pt}
\setlength{\emergencystretch}{3em}
\newcommand{\N}{\mathbb{N}}
\newcommand{\Z}{\mathbb{Z}}
\newcommand{\Q}{\mathbb{Q}}
\newcommand{\R}{\mathbb{R}}
\newcommand{\C}{\mathbb{C}}
\newcommand{\F}{\mathbb{F}}
\newcommand{\A}{\mathbb{A}}
\newcommand{\PP}{\mathbb{P}}
\newcommand{\E}{\mathbb{E}}
\newcommand{\eps}{\varepsilon}
\newcommand{\dd}{\,\mathrm{d}}
\begin{document}
\section*{UnsolvedMath ID: 701020}
\textbf{Source problem code:} catalogue entry
\subsection{Definitions and mathematical statement}
A group \(G\) is finitely presented if there exist nonnegative integers \(d,s\) and words \(r_1,\ldots,r_s\) in the free group on \(x_1,\ldots,x_d\) such that
\[
G\cong\langle x_1,\ldots,x_d\mid r_1,\ldots,r_s\rangle.
\]
The exponent of \(G\), when it exists, is the least positive integer \(n\) such that \(g^n=1_G\) for every \(g\in G\). Thus a nontrivial group has exponent equal to a prime \(p\) precisely when every element satisfies \(g^p=1_G\).

\textbf{Question (N. D. Gupta).} For every prime \(p\) and every finitely presented group \(G\) of exponent \(p\), must \(G\) be finite? Equivalently, is the implication
\[
\left(G\text{ is finitely presented and }(\forall g\in G)\;g^p=1_G\right)
\Longrightarrow |G|<\infty
\]
valid for every prime \(p\) and every group \(G\)? The equivalent formulation also includes the trivial group, which is finite.
\subsection{Short English statement}
Does every finitely presented group whose exponent is a prime have finitely many elements? The question ranges over all primes.
\subsection{Sources}
\begin{itemize}
\item Local UnsolvedMath registry, supplied complete source record: permanent ID \(701020\), exact source problem code \texttt{catalogue entry}, title \emph{Kourovka Notebook Problem 14.35}.
\item The supplied record attributes Problem 14.35 to N. D. Gupta, proposed in 1999, and identifies its source as \emph{The Kourovka Notebook (Unsolved Problems in Group Theory)}, 14th issue (1999), with the 21st issue edition (2026), arXiv:1401.0300v48. The collection description names E. I. Khukhro and V. D. Mazurov as editors.
\item The local record labels the problem \emph{open}. This is supplied source metadata, not an independently verified assessment of the literature. No external source was consulted for this draft.
\end{itemize}
\end{document}
